\documentclass[11pt]{article}
\usepackage[a4paper,margin=1in]{geometry}
\usepackage[utf8]{inputenc}
\usepackage[T1]{fontenc}
\usepackage{amsmath,amssymb}
\usepackage{hyperref}
\usepackage{enumitem}
\usepackage{booktabs}

\title{\textbf{Recombination in a Self-Modifying Network Search:\\
A Building-Block Validation and a Structural Negative Result}}
\author{Vasili Gavrilov\thanks{Independent researcher. Physics and computer
science by education; a professional software engineer with a background in
numerical optimization and machine learning, including rule engines and expert
systems. Reference implementation: \url{https://github.com/Anode1/SMBPANN}.}}
\date{2026}

\begin{document}
\maketitle

\begin{abstract}
Ordinarily a neural network must be handed its shape before it can learn: how many
layers, how wide, which activation. The system studied here is handed only the
training data and a target error, and finds the rest itself by evolution: its depth
and width, whether each layer is dense or a weight-sharing convolutional one, and its
own learning rate, momentum, and activation. Onto that mutation-driven search we add
\emph{crossover}, recombination of two parents, following the intuition that mixing
building blocks should reach a good architecture faster than mutating one lineage.
We report an honest pair of results. First, on a problem with genuine separable
building blocks, concatenated deceptive trap functions, our one-point crossover
beats mutation decisively, and the advantage \emph{grows} with the number of blocks
(from $1.8\times$ fewer generations at four blocks to mutation failing outright while
crossover still succeeds at sixteen), the textbook fingerprint of a real
recombination effect. Second, on the neural architecture search, crossover is
on-par with mutation, no better. We diagnose why, with measurements: on our task the
error is set far more by the training hyper-parameters than by the architecture, and
architecture matters only through a brittle interaction with them; and because each
candidate is trained from scratch, the genome carries architecture but not the
learned solution, so there is nothing for recombination to transfer. Neural capacity
is fungible and its solutions live in weights, not in the architecture genome, which
is the structural reason random search is a strong architecture-search baseline and
crossover is routinely dropped. We identify \emph{weight inheritance}, carrying a
child's trained parts across generations so a useful combination keeps evolving, as
the change that would give recombination something to transfer, and mark it as the
next step.
\end{abstract}

\section{Introduction}

Neural architecture search (NAS) automates a choice a human normally makes: the shape
of the network. Its methods differ mainly in the search strategy: reinforcement
learning \cite{zoph2017}, gradient-based relaxation \cite{darts2019}, Bayesian
optimization, random search, and evolutionary search \cite{elsken2019}. Ours is
evolutionary: the topology is evolved while each candidate's weights are trained by
back-propagation \cite{real2019}. Ordinarily the architecture and training knobs are
arguments the user supplies, \textbf{predict(topology, activation, \ldots)}; the
system here is handed only the data and a target error, \textbf{self\_modifying\_predict(data,
error)}, and discovers the rest. The engine grew out of the author's 1997 seminar
thesis, which derived the general multi-layer delta rule the training still uses
\cite{thesis1997}.

A recurring, humbling result in this field is that \emph{random search is hard to
beat} \cite{bergstra2012,li2019}. This paper asks whether \emph{recombination} can
do better than mutation at the architecture search: whether crossing two good
parents, mixing their building blocks, reaches a target error faster than mutating
one. The self-adaptive machinery and the evolution-strategy view we lean on are
standard in the optimization literature the author has worked in \cite{rechenberg1973}.
The honest answer we reach is two-sided, and the second side is a structural fact
about neural architecture search that a positive control makes precise.

\section{Method}

\subsection{The self-modifying search}

A network is a feed-forward stack of dense and convolution-like layers. For a dense
layer $\ell$ the forward pass is $a^{(\ell)} = g(W^{(\ell)}a^{(\ell-1)} + b^{(\ell)})$,
with $g$ a per-network hidden activation (logistic sigmoid, $\tanh$, or ReLU) chosen
by the genome; the output layer is a dense sigmoid. A convolutional layer applies $F$
shared kernels of width $K$ across the previous layer, so its width is derived,
$d_\ell = F(d_{\ell-1}-K+1)$, and its weight count is only $F\,K$: the local,
weight-sharing structure LeCun introduced for images \cite{lecun1989}, here a move
the search can take. Weights are trained by back-propagation with momentum (the
generalized multi-layer delta rule of the 1997 thesis \cite{thesis1997}), initialized
uniformly in $[-2.4/f_i,\,2.4/f_i]$ with $f_i$ the fan-in. All arithmetic is 32-bit
float; all randomness flows through a seeded xorshift generator.

A \emph{genome} carries the topology (per hidden layer: dense or convolutional, and
its free parameters) and its own training hyper-parameters (learning rate, momentum,
activation), plus a self-adaptive mutation rate. One \emph{mutation} makes a single
small change: perturb a layer, flip a layer dense$\leftrightarrow$conv, or add or
remove a layer, within depth and width bounds. Reproduction is self-adaptive in the
sense of evolution strategies \cite{rechenberg1973}: an offspring inherits and
log-normally perturbs its parent's mutation rate, applies that many mutations, and
co-evolves the hyper-parameters. The search keeps a population, retains the best few
(elitism), and refills by reproduction, against a matched random-search control that
draws the same number of fresh random genomes each generation.

\subsection{The crossover operator}

To the mutation-only search we add optional \emph{crossover} (\texttt{evolve -X}, the
percent of offspring made by recombination). A crossover child is built from two
elite parents by a one-point splice of their hidden-layer lists, a prefix of one and
a suffix of the other, so it inherits building blocks and a new depth from both; its
training hyper-parameters are a blend (geometric mean of the learning and mutation
rates, arithmetic mean of momentum, one parent's activation). The child is then
mutated as usual. With \texttt{-X 0} the operator is inert and the run is
byte-identical to the mutation-only search.

\section{Experiments}

\subsection{Does the operator work? A building-block validation}

Before asking whether crossover helps the neural search, we check that the operator
does what crossover is supposed to do, on a problem built to have separable building
blocks. We use concatenated deceptive trap functions: a genome of $M$ blocks of $K$
bits, where a block scores $K$ if all its bits are one, else $K-1$ minus its number
of ones, a deceptive gradient luring a bit-flipping search toward all-zero. The
global optimum (all ones) can be reached only by assembling correctly-solved blocks,
which is exactly what crossover recombines and mutation cannot climb to. We race the
same GA (tournament selection, elitism) with one-point crossover against mutation-only,
over 200 seeds, and scan the number of blocks (\texttt{make bbtest}).

\begin{table}[t]
\centering
\caption{Crossover versus mutation on concatenated deceptive traps ($K=4$-bit blocks,
population 1000, tournament 3, 200 seeds). As the number of building blocks grows,
crossover's advantage grows: mutation's cost explodes and it begins to fail outright,
while crossover keeps solving. This is the fingerprint of a genuine building-block
effect.}
\label{tab:bb}
\begin{tabular}{lcccc}
\toprule
blocks $M$ (genome) & mutation-only & crossover $+$ mutation & speedup \\
\midrule
4  (16 bits) & 9.2 gens, 200/200  & 5.1 gens, 200/200  & $1.8\times$ \\
8  (32 bits) & 70.8 gens, 200/200 & 14.2 gens, 200/200 & $5\times$ \\
12 (48 bits) & 257 gens, 185/200  & 22.6 gens, 200/200 & $11\times$, more reliable \\
16 (64 bits) & 363 gens, \textbf{28/200} & 42.3 gens, \textbf{199/200} & mutation fails \\
\bottomrule
\end{tabular}
\end{table}

Table~\ref{tab:bb} is decisive. At sixteen blocks mutation-only reaches the global
optimum on only $28$ of $200$ seeds; crossover reaches it on $199$, in a fraction of
the generations. The operator and the methodology are sound: when separable building
blocks exist, crossover finds and combines them, and its advantage scales with the
amount of structure.

\subsection{Does crossover help the neural search? No}

We then run the same on/off comparison on the architecture search: for each of a set
of seeds, a fresh synthetic classification task ($d=8$ inputs, a random-projection
sine label), the self-modifying search run twice from the identical initial
population, once mutation-only and once with crossover (every offspring recombining
two elites), each until its best validation error reaches a target, or a generation
cap. Over the paired runs collected, crossover and mutation are indistinguishable:
each reached the target on the same fraction of seeds, with the same median
generations-to-target and an even head-to-head. Adding recombination did nothing.
The next subsection explains why, and the explanation is not that the operator is
weak (Table~\ref{tab:bb}), but that the neural search has no building blocks for it
to recombine.

\subsection{Why: architecture barely matters, and the genome carries no solution}

Two measurements on a fixed task tell the story. First, architecture alone does not
solve the task at all: trained with default hyper-parameters, every topology we tried,
dense and convolutional, wide and deep, scored between $0.25$ and $0.59$ test-MSE, at
or worse than the $0.25$ a constant predictor scores. Second, the \emph{hyper-parameters}
do the work, and brittly: holding the topology at a small $8,4,1$ and varying only
the learning rate, momentum, and activation swings the error from $0.09$ to $0.65$;
and with a good hyper-parameter set fixed, varying the topology is non-monotone and
narrow, one width learning ($8,4,1 \to 0.09$) while its neighbours fail
($8,8,1 \to 0.37$, $8,16,1 \to 0.41$). Each architecture wants its own learning rate;
evaluated with the wrong one, a fine architecture looks terrible. The landscape is
rugged and the two dimensions are entangled, the worst case for a directed search,
and crossover of a topology block breaks the co-adaptation that made the block good.

A separate probe removes even that entanglement and still finds no building blocks. A
parallel-branch network (\texttt{make modnas}, back-propagation gradient-checked
against finite differences, $0/40$ mismatches) sums $B$ independent branches into a
multi-output prediction, on an additively separable task where each output depends on
its own group of inputs, the closest neural analogue of the trap's independent blocks.
Yet the branch count barely changes the error, and no configuration solves the task
cleanly, because \emph{neural capacity is fungible}: a single sufficiently-wide branch
computes what several specialised ones would, so the architecture does not decompose
into independent modules that recombination could assemble.

Underlying both is one fact. In the trap problem the bits \emph{were} the solution, so
crossover transferred solutions. In architecture search the genome is architecture and
the solution is \emph{weights}, retrained from scratch for every candidate, so the
genome carries no solution to transfer, and architectural choices are, per above,
fungible and interchangeable. There is nothing for recombination to recombine.

\subsection{A note on mutation versus random}

For completeness: the mutation-only self-modifying search does edge its matched random
control on time-to-target, by a small statistically-clear margin over many seeds. But
the diagnosis above reframes that edge: since the architecture barely moves the error,
the advantage comes from co-evolving the \emph{hyper-parameters} (the self-adaptive
learning rate, momentum, activation), not from searching architecture. It is a
hyper-parameter result wearing an architecture-search costume, and we do not lean on
it.

\section{Discussion}

The two results compose into one honest statement. Crossover's building-block
advantage is real and provable, and it scales with structure (Table~\ref{tab:bb}). But
neural architecture search, as usually posed, \emph{structurally lacks that structure},
for two reasons we have measured: neural capacity is fungible, so architectures do not
decompose into independent modules; and solutions live in the trained weights, which a
retrain-from-scratch search does not inherit, so the genome carries nothing to
transfer. Together these are why random search is a strong NAS baseline and why the
field \cite{real2019} routinely drops crossover. Our contribution is not a new method
but a clean demonstration, with a positive control, of exactly when recombination
helps and why it usually cannot here.

\paragraph{Limitations.} The scale is small and the task synthetic; the neural
crossover comparison is preliminary (few paired seeds), though its flatness is
consistent with the mechanism the fixed-task measurements establish. Mean squared
error stands in for classification error. None of this is hidden.

\section{Conclusion and future work}

Handed only the problem and a target error, the search discovers architecture, layer
type, and training hyper-parameters, and stops when the error is met. Adding
recombination did not speed it up, and we now understand structurally why: there are
no building blocks in a fungible-capacity, weights-not-inherited architecture genome.
The natural fix, and the next step, is \emph{weight inheritance}: let a child carry
its parents' trained parts rather than retrain from zero, so a useful combination that
evolution has found keeps evolving instead of restarting, which should both give
recombination a solution to transfer and converge faster. That turns the genome from
architecture-only into architecture-plus-learned-state, the form in which the
building-block intuition, validated in Table~\ref{tab:bb}, can finally apply to
networks. Two further steps remain independent of it: a two-dimensional convolution
and an image task, where architecture is decisive rather than fungible; and a bridge
to a standard tabular NAS benchmark.

\section*{Reproducibility}

Every result is regenerated from \cite{repo} with a single command. The engine builds
with \texttt{make} (C99, no dependencies) into \texttt{smbpann} (worker),
\texttt{evolve} (search and control, with \texttt{-X} for crossover), and
\texttt{gentask} (task generator); \texttt{make ut} runs the unit and gradient-check
suite under AddressSanitizer and UBSan. The building-block validation
(Table~\ref{tab:bb}) is \texttt{make bbtest} then \texttt{./bbtest M pop}; the
parallel-branch probe is \texttt{make modnas}. The neural crossover A/B is
\texttt{scripts/xover\_ab.sh}; the search and task are controlled by environment
variables (\texttt{DIM}, \texttt{FREQ}, \texttt{TARGET}, \texttt{POP}, \texttt{ELITE},
\texttt{XPCT}, and so on). Each run uses a fixed 32-bit xorshift generator, so the
numbers are identical on re-running.

\begin{thebibliography}{99}
\small
\bibitem{rumelhart1986} D.~E. Rumelhart, G.~E. Hinton, R.~J. Williams.
  Learning representations by back-propagating errors. \emph{Nature} 323, 1986.
\bibitem{thesis1997} V.~Gavrilov. \emph{Backpropagation Feed-Forward Neural
  Networks}. Seminar in Artificial Intelligence, Open University of Israel, 1997.
\bibitem{rechenberg1973} I.~Rechenberg. \emph{Evolutionsstrategie}. 1973;
  H.-P. Schwefel, \emph{Numerical Optimization of Computer Models}, 1981.
\bibitem{holland1975} J.~H. Holland. \emph{Adaptation in Natural and Artificial
  Systems}. University of Michigan Press, 1975 (schemata and building blocks).
\bibitem{deb1993} K.~Deb, D.~E. Goldberg. Analyzing Deception in Trap Functions.
  \emph{Foundations of Genetic Algorithms} 2, 1993.
\bibitem{lecun1989} Y.~LeCun, B.~Boser, J.~S. Denker, et al. Backpropagation Applied
  to Handwritten Zip Code Recognition. \emph{Neural Computation} 1(4), 1989.
\bibitem{bergstra2012} J.~Bergstra, Y.~Bengio. Random Search for Hyper-Parameter
  Optimization. \emph{JMLR} 13, 2012.
\bibitem{zoph2017} B.~Zoph, Q.~V. Le. Neural Architecture Search with Reinforcement
  Learning. \emph{ICLR} 2017.
\bibitem{real2019} E.~Real, A.~Aggarwal, Y.~Huang, Q.~V. Le. Regularized Evolution
  for Image Classifier Architecture Search. \emph{AAAI} 2019.
\bibitem{darts2019} H.~Liu, K.~Simonyan, Y.~Yang. DARTS: Differentiable Architecture
  Search. \emph{ICLR} 2019.
\bibitem{li2019} L.~Li, A.~Talwalkar. Random Search and Reproducibility for Neural
  Architecture Search. \emph{UAI} 2019.
\bibitem{elsken2019} T.~Elsken, J.~H. Metzen, F.~Hutter. Neural Architecture Search:
  A Survey. \emph{JMLR} 20, 2019.
\bibitem{repo} V.~Gavrilov. SMBPANN: a self-modifying backpropagation feed-forward
  neural network. Source code, \url{https://github.com/Anode1/SMBPANN}.
\end{thebibliography}

\end{document}
